\section{Priced local families give the three face estimates}
\label{sec:ex-annulus-fleet-faces}

We derive the upper-face bound, the full lower-face bound, and the
lower-face-minus-halo bound used in the annulus estimate. The proof is the
local price calculation in the selected construction. We then verify the
scale and data choices for the two populations used by the annulus: all
centered configurations and the regulator-thin shell. We derive the
selected row's contact identities and the scale-weighted face prices,
including the separate shell restriction required by the unregulated
reference inverses.

\subsection{Derivative prices}

For analytic matrix paths $M_i(t)$ near $t=0$, with values in a fixed
finite-dimensional matrix space of dimension $d$, write
\[
 \mathsf P_{r}(M;C)\quad\Longleftrightarrow\quad
 M_i\text{ is analytic near }0\text{ and }
 r_i^j\,\|M_i^{(j)}(0)\|_{\mathrm{op}}\leq C
 \quad(j=0,1,2),
 \tag{F1}
\]
uniformly in the index $i$. Here $\|\cdot\|_{\mathrm{op}}$ is the
Euclidean operator norm; it is submultiplicative, and
$|\operatorname{tr}M|\leq d\|M\|_{\mathrm{op}}$, where $d$ is the matrix
dimension. Constants may be enlarged by a fixed factor without changing
the notation.

The elementary rules used below are as follows. If $M,N$ have prices
$(r,C)$ and $(r,D)$, then $M+N$ has price $(r,C+D)$ and $MN$ has price
$(r,4CD)$. For the product, the zeroth and first derivatives are the
ordinary product rule; the second derivative has four terms, each bounded
by submultiplicativity and the corresponding two price bounds. If $a>0$
and $C,D\geq0$, $M(0)$ is invertible and
$\|M(0)^{-1}\|_{\mathrm{op}}\leq D/a$, while $M$ has price $(r,aC)$,
differentiating $M^{-1}M=I$ once and twice gives
\[
 \mathsf P_r(M^{-1};\mathcal I(C,D)/a),\qquad
 \mathcal I(C,D)=D+2CD^2+2C^2D^3.
 \tag{F2}
\]
Indeed, $(M^{-1})'=-M^{-1}M'M^{-1}$ and
$(M^{-1})''=2M^{-1}M'M^{-1}M'M^{-1}-M^{-1}M''M^{-1}$;
submultiplicativity bounds the terms by the three summands in (F2).
Finally,
\[
 \mathsf P_r(M;C r_i^{-2})
 \ \Longrightarrow
 |(\operatorname{tr}M_i)''(0)|\leq dC r_i^{-4},
 \tag{F3}
\]
by the $j=2$ case of (F1) and the trace inequality. These are the
finite-dimensional price rules underlying the estimates below.

\subsection{Upper, lower, and halo faces}

For each configuration $i$, let $x_i$ be its background, $e_i$ its
boost, $\mu_i>0$ its regulator and $T_i>0$ its scale. The local geometry
used here is
\[
 \|e_i\|\leq1,\qquad \|x_i\|\leq T_i\leq r_*,\qquad
 0<\mu_i\leq T_i.
 \tag{F4}
\]
Let $S_i$ denote the total stationary denominator (interacting plus
reference), and let $J_i$ and $J_i^{\rm ref}$ be the two two-leg arch
insertions. Write $K_{0,i}=S_i(0)^{-1}$ for the inverse at the
unshifted background and $K_{1,i}(t)$ for the inverse at the displaced
background $x_i+t e_i$. Write $D_{0,i}$ and $D_{1,i}(t)$ for the
corresponding free/reference inverses. The unshifted inverse is constant
in $t$; the shifted inverse varies. With $B_i^-,B_i^+$ the displaced
insertion legs, define the scalar upper and lower faces by
\[
 U_i(t)=\operatorname{tr}\left[K_{0,i}(J_i(t)+J_i^{\rm ref}(t))\right],\qquad
 L_i(t)=\operatorname{tr}\left[K_{0,i}B_i^-(t)K_{1,i}(t)B_i^+(t)\right].
 \tag{F5}
\]
The halo is $H_i(t)=\operatorname{tr}[D_{0,i}\widetilde B_i^-(t)
D_{1,i}(t)\widetilde B_i^+(t)]$, with the four factors the
corresponding free/reference values. Thus it replaces both inverses and
both legs.

The selected row transfer yields the following prices, with one constant
$C_0\geq0$ independent of $i$:
\begin{align}
 &\mathsf P_{T_i}(K_{0,i};C_0T_i^{-2}),\quad
   \mathsf P_{T_i}(K_{1,i};C_0T_i^{-2}), \label{F6a}\\
 &\mathsf P_{T_i}(J_i^{\rm ref};C_0\mu_i^2T_i^{-2}),
 \qquad \mathsf P_{1}(J_i;C_0), \label{F6b}\\
 &\mathsf P_{T_i}(B_i^\pm;C_0T_i), \label{F6c}\\
 &\mathsf P_{T_i}(K_{0,i}-D_{0,i};C_0\,\delta_iT_i^{-2}),\quad
   \mathsf P_{T_i}(D_{0,i};C_0T_i^{-2}),\notag\\
 &
   \mathsf P_{T_i}(D_{1,i};C_0T_i^{-2}),\quad
   \mathsf P_{T_i}(K_{1,i}-D_{1,i};
                    C_0\,\delta_iT_i^{-2}),\notag\\
 &\mathsf P_{T_i}(B_i^\pm-\widetilde B_i^\pm;
                      C_0\,\delta_iT_i),\label{F6d}
\end{align}
where $C_0\geq0$ and $\delta_i=T_i+\mu_i^2/T_i^2$. The tilded factors are the
reference/halo factors just described. Also
$\|K_{0,i}\|_{\mathrm{op}}\leq C_0T_i^{-2}$. The ordinary face prices
(F6a)--(F6c) apply to the all-centered population. The extra differences
(F6d) are required only on the shell $\mu_i\leq T_i=\|x_i\|_2$; at the
center $x_i=0$ the unregulated halo inverse need not exist.
The derivation from this same selected row and the reference field is given
below; in particular, the free inverses in (F6d) are used only on the shell.

\begin{proposition}[Face estimates]
\label{prop:ex-fleet-face-bounds}
Under (F4)--(F6c), there are constants $A,B,C\geq0$, independent of
$i$, such that
\[
 |U_i''(0)|\leq \frac{A}{T_i^2}
              +B\frac{\mu_i^2}{T_i^6},\qquad
 |L_i''(0)|\leq\frac{C}{T_i^4},
 \tag{F7}
\]
and, whenever the additional shell prices (F6d) hold,
\[
 |L_i''(0)-H_i''(0)|\leq C\frac{\delta_i}{T_i^4}.
 \tag{F8}
\]
\end{proposition}
\begin{proof}
For the reference part of the upper face, (F6a)--(F6b) and the product
rule give price $C\mu_i^2/T_i^4$ at scale $T_i$. By (F3), its second
derivative is at most $C\mu_i^2/T_i^6$. For the interacting part,
$S_i(0)^{-1}$ is constant in $t$ and has norm at most $C_0/T_i^2$;
combine this with the unit-scale price of $J_i$ and apply the trace
bound at unit scale. This gives $A/T_i^2$. Adding proves the first
inequality in (F7).

The lower-face package supplies price $C/T_i^2$ at scale $T_i$ for the
four-factor product in (F5): each inverse has price proportional to
$T_i^{-2}$ and each displaced leg has price proportional to $T_i$.
Applying (F3) yields $C/T_i^4$, proving the second inequality.

For the difference estimate, assume in addition (F6d), as is valid on the
shell. Since $T_i\leq r_*$ and $\mu_i\leq T_i$, we have
$\delta_i\leq r_*+1$. The sum rule for prices and (F6c)--(F6d) therefore
give each tilded leg the ordinary price
$\mathsf P_{T_i}(\widetilde B_i^\pm;C_1T_i)$, with $C_1$ depending only
on $C_0$ and $r_*$. Use the exact ordered telescoping identity
\[
 \begin{aligned}
 A_1A_2A_3A_4-B_1B_2B_3B_4
 &= (A_1-B_1)A_2A_3A_4+B_1(A_2-B_2)A_3A_4\\
 &\quad+B_1B_2(A_3-B_3)A_4+B_1B_2B_3(A_4-B_4).
 \end{aligned}
 \tag{F9}
\]
Both inverse differences, at the unshifted and shifted legs, have price
$C\delta_i/T_i^2$ by (F6d); the leg differences have price
$C\delta_iT_i$ by (F6d). The remaining inverse factors cost $C/T_i^2$ and the remaining legs
cost $CT_i$. In each of the four summands the product rule therefore
gives price at most $C\delta_i/T_i^2$ at scale $T_i$. Apply (F3) to
the trace of their sum. Its second derivative is bounded by
$C\delta_i/T_i^4$, as claimed.
\end{proof}

\subsection{The Hstar populations and the two scale choices}

Let $G$ be simple and fix the reference family, positive regulator fee,
grey bound, and local row data. An Orion datum supplies a centered track
$x$, regulator $\mu$, color $c$, bed, clip and a local family $l$;
the row condition is imposed on the combined family $l+V$. An Hstar
datum is such an Orion datum together with the condition that $x$ lies in
the closed centered box. The selected population uses the axial mark
$(1,c)$ and the unit axial boost $e_0$.

For each selected row, the path bounds give caps for the tab, the
one-leg insertions, and the two-leg insertion, together with the two
soft contacts. The regulated reference path supplies its own bounds.
Combining the paths gives the total stationary denominator and its
scale-$T^2$ lower bound. The derivation of (F6) below shows explicitly
how these data yield the required prices. In particular, no
scale-weighted estimate is inferred from analyticity alone: the rates
come from the path budget, its one-leg weight $T^{-1}$, and its two-leg
weight $T^{-2}$; the extra shell estimates use subtraction identities.

For a centered-box track, each coordinate has modulus at most $\pi$,
so $\|x\|_2\leq\sqrt{4\pi^2}=2\pi$; the regulator satisfies
$0<\mu\leq1$. The
all-centered scale is
\[
 T=\|x\|_2+\mu,
 \qquad 0<T\leq1+2\pi=:r_*.
 \tag{F10}
\]
Thus (F4) holds and (F7)'s upper estimate gives the first two-wave
bound after combining it with the lower-face estimate. On the shell
$\mu\leq\|x\|_2$, reset the scale to
$T=\|x\|_2$. Then $0<\mu\leq T\leq r_*$, and the same package gives
both the upper estimate and (F8). Use the low-face identity of
Proposition~\ref{prop:ex-two-wave} to replace the halo by the continuum
pole. Since $T\leq r_*$,
\[
 T^{-2}\leq r_*T^{-3},\qquad T^{-2}\leq r_*^2T^{-4},\qquad
 \frac{\mu^2}{T^6}\leq\frac1{T^4},\qquad
 \frac{\delta}{T^4}=\frac1{T^3}+\frac{\mu^2}{T^6}.
\]
Indeed, the response is half the difference of interacting and reference
upper-minus-lower traces. The upper estimate and lower-face estimate in
(F7), with $\mu\leq T$, therefore give $|\kappa(x)|\leq aT^{-4}$.
For the shell discrepancy, the exact algebra is
\[
 \kappa(x)-\Pi_G(x)
 =\frac12\left(\frac12U_I''-U_R''
                  +(L_R''-H_R'')-\frac12(L_I''-H_I'')\right),
 \tag{F12}
\]
where the reference/interacting halo combination is the pole identity
in Proposition~\ref{prop:ex-two-wave}. The two upper terms are bounded by
$C/T^2+C\mu^2/T^6$, and the lower remainders by $C\delta/T^4$;
the displayed scale inequalities give the two powers in (F11).
Enlarge constants once to cover the bounded part of the centered box
outside the unit ball. We obtain constants $a,b$ independent of $x$ and
$\mu$ such that
\[
 |\kappa(x)|\leq a(\|x\|_2+\mu)^{-4},\qquad
 |\kappa(x)-\Pi_G(x)|\leq b\|x\|_2^{-3}
      +b\mu^2\|x\|_2^{-6}\quad(\mu\leq\|x\|_2\leq1).
 \tag{F11}
\]
These are the uniform local estimates consumed in
Section~\ref{sec:ex-annulus-two-wave}.

\subsection{The local row bridge to the interacting family}

Here are the exact row identities behind the local contact conditions.
For a selected bed $B_i$, write $M_i(p)=B_i.\mathrm{metric}(p)$ and let
$M_*(p)$ be the constant leader metric. Let $D(p)$ be the first edge
symbol, with components $q_\nu(p)$, and let $A_0(p)$ be the zero-form
symbol. The color-diagonal projection of the direction matrix is denoted
by $\mathsf T$. Put
\[
 G_i(p)=\text{the two-point row of the interacting film},\qquad
 J_{i,b}(p,k)=\text{its marked two-leg row}.
\]
For the selected construction, the row package has a common analytic
strip and fee bound; it satisfies the denominator matching identity
\[
  (\text{interacting film}+\text{regulator film})_{\rm tab}(p)
       =\text{the bed denominator}(p),
\tag{F13}
\]
the all-grade floor (in particular, the grade-one normalization
$M_i(0)=M_*(0)$), the local action soft order-three condition, and the
vacuum matching condition (the color-diagonal vacuum row equals the
finite Wilson row). These properties all concern the same selected bed
and the same combined local row. The regulator satisfies $0<\mu\leq1$
and its tab is the scalar multiple $s_\mu(p)I$ on the strip. They are
properties of the selected row package, rather than consequences of
analyticity alone.

Let $\omega(p)=\sum_\nu q_\nu(-p)q_\nu(p)$. The exact row and free
identities use the color-diagonal extension
$\mathsf T(A)_{(i,c),(j,d)}=\mathbf 1_{c=d}A_{ij}$ of a direction
matrix; this extension preserves matrix sums and products. They are
\[
 G_i(p)-|p|^2I
   =\mathsf T\!\left[P_1(-p)^T(M_i(p)-M_*(p))P_1(p)\right]
       +(\omega(p)-|p|^2)I,
\tag{F14}
\]
\[
 J_{i,b}(p,k)-v_b(p,k)
   =J_{i,b}^{\rm loc}(p,k)
       +\bigl(J_{i,b}^{\rm vac}(p,k)-v_b(p,k)\bigr),
\tag{F15}
\]
and
\[
\mathsf T\!\left[P_1(-p)^TM_*(p)P_1(p)
                +P_0(p)P_0(-p)^T\right]=\omega(p)I.
\tag{F16}
\]
Here $v_b$ is the linear vector insertion determined by the first radial
coefficient of the vacuum marked row.

\begin{proof}[Derivation of (F14)--(F16)]
Write $P_0(p)$ and $P_1(p)$ for the zero-form and first-edge symbols, so
$A_0=P_0$ and $D=P_1$. The selected denominator matching equates the
tab of the interacting film plus its regulator film to the prescribed
bed denominator. The two sides expand as
\[
 \begin{split}
 (\mathrm{interacting}+\mathrm{regulator})_{\rm tab}(p)
   &=G_i(p)+s_\mu(p)I,\\
 (\mathrm{bed\ denominator})(p)
   &=\mathsf T\!\left[P_1(-p)^TM_i(p)P_1(p)
           +P_0(p)P_0(-p)^T\right]+s_\mu(p)I.
 \end{split}
\]
The regulator tab is scalar on the slot space, so equality cancels that
term entrywise and gives, on real momenta,
\[
 G_i(p)=\mathsf T\!\left[P_1(-p)^TM_i(p)P_1(p)
                    +P_0(p)P_0(-p)^T\right].
\]
The first summand is the metric cross product; the second is the zero-form
product. The cancellation is exact, so it introduces no regulator
remainder. For the leader bed, the degree-zero exterior identity gives
\[
 \mathsf T\!\left[P_1(-p)^TM_*(p)P_1(p)
              +P_0(p)P_0(-p)^T\right]=\omega(p)I.
\]
Indeed, expand the first product by direction. The longitudinal rank-one
part of the zero-form product cancels the complementary part of the
first-edge product, leaving $\sum_\nu q_\nu(-p)q_\nu(p)$ times the
identity. The exact degree-zero exterior identity is
$P_0(p)P_0(-p)^T+P_1(-p)^TM_*P_1(p)=\omega(p)I$,
since $\omega(p)=\sum_\nu q_\nu(-p)q_\nu(p)$ by definition. Since
$M_*$ is the constant identity metric, this proves
(F16). Subtracting the
leader identity from the row-to-film identity proves (F14), with the
lattice-dispersion term $\omega(p)-|p|^2$ retained explicitly.

For the marked row, the local film is the sum of the selected local
action and the fixed free quadratic film. Linearity of the marked
coefficient gives (F15). The soft-3 hypothesis assigns radial order two
to the one-mark layer; composing that layer with the linear two-momentum
map gives the required second-order vanishing. The vacuum term is identified by grey
matching with the Wilson three-leg coefficient. Write $\tau(u)_a=iu_a$
and denote the color structure coefficient by $c_{bxy}$. Differentiating
the three plaquette placements gives its first radial coefficient
\[
 [-(2\tau(k)_a+\tau(p)_a)\delta_{ij}
   +(\tau(k)_j+2\tau(p)_j)\delta_{ia}
   +(\tau(k)_i-\tau(p)_i)\delta_{ja}]c_{bxy}.
\]
Each placement contributes one of the terms in the defining plaquette
coefficient; summing the two orientations and using
$\tau(k+p)=\tau(k)+\tau(p)$ gives this formula. The fixed free-film
marked coefficient has first radial coefficient
\[
 [-\tau(k)_j\delta_{ia}-\tau(k+p)_i\delta_{ja}]c_{bxy}.
\]
Adding these expressions and simplifying gives
\[
 v_{b,ij}(p,k)=
 [-(2\tau(k)_a+\tau(p)_a)\delta_{ij}
      +2\tau(p)_j\delta_{ia}-2\tau(p)_i\delta_{ja}]c_{bxy}.
\]
This is homogeneous of degree one. Thus the vacuum marked row minus
$v_b$ has zero value and zero first radial coefficient; analyticity
implies that it has radial order at least two. With the soft local term,
(F15) has the required quadratic contact. The common strip bounds then
give uniform compact derivative bounds by Cauchy's formula.
\end{proof}

\begin{lemma}[Row data give the priced insertion family]
\label{lem:ex-row-to-floe}
For a bounded family of selected rows with the common strip/fee bounds,
denominator matching, all-grade floor, soft order-three local action and grey
vacuum matching just described, the interacting film has uniform
compact derivative prices through order three. Its two-point row has a
cubic radial remainder after subtraction of $|p|^2I$, and its marked
row has a quadratic radial remainder after subtraction of $v_b$.
Consequently it supplies the compact caps and the two contact orders
used in the ordinary (non-halo) part of (F6b)--(F6c).
\end{lemma}
\begin{proof}
For each row layer and each fixed pair of matrix entries, the common
strip fee and Cauchy's estimate bound derivatives of every fixed order
on a compact real ball. The tab, spot and arch are fixed finite linear
combinations of these entries after fixed linear maps of momentum; the
chain rule therefore gives uniform compact derivative caps through the
orders used here. A finite sum of these caps bounds the matrix operator
norm, since the matrix dimension is fixed.

For the two-point contact, (F14) and the grade-one floor give
$M_i(0)-M_*(0)=0$. Analyticity makes this metric difference $O(|p|)$.
Each edge symbol vanishes linearly, $q_\nu(p)=ip_\nu+O(|p|^2)$, so the
metric term in (F14) is $O(|p|^3)$. Expanding each product
$q_\nu(-p)q_\nu(p)$ shows $\omega(p)-|p|^2=O(|p|^3)$ as well.

For the marked row, the local soft condition contributes $O((|p|+|k|)^2)$
and the vacuum-minus-$v_b$ term has the same order by the preceding
first-coefficient calculation. Thus (F15) has quadratic contact.
To convert these radial orders into derivative prices, note that if an
analytic scalar entry $g$ has radial order at least $m$, then each
homogeneous Taylor polynomial of degree $j<m$ vanishes on every real
direction and hence vanishes identically. Taylor's formula with integral
remainder applied to $D^n g$ on the segment from $0$ to the base point
$a$ gives
\[
 \|D^n g(a)\|\leq K\|a\|^{m-n}/(m-n)!,\qquad n<m,
\]
when the $m$th derivative is bounded by $K$ on that segment. For
$\|a\|\leq T$, this gives $T^n\|D^ng(a)\|\leq KT^m$. At orders
above the contact order use the compact derivative cap and the uniform
upper bound $T\leq r_*$. This proves the local compact and contact
prices. The trajectory-price argument below derives the total-field
prices in (F6a) and the shell-separation prices in (F6d) from these
contacts and the reference path budgets.
\end{proof}

\paragraph{Hstar to the local row package.}
For every Hstar datum the associated Orion datum contains one selected
bed, clip, regulator, and local action. Its denominator matching field is
the identity
(F13), its fee field supplies the uniform analytic strip, and its
all-grade floor supplies the grade-one normalization. Its local soft
field is soft order three. Applying this row argument to the single
combined family $l+V$ gives the interacting row used by the annulus.
The common free reference is the quadratic film. If the external vacuum
$V$ is grey and has the required fees, the finite-volume Wilson matching
above identifies its marked vertex and hence its linear term. These
statements establish the row contacts from the selected Hstar data. The
scalar regulator term and its positive floor are then included in the
total denominator; the next price calculation uses the free comparison
only on the shell, where its inverse exists.

Assume $\epsilon>0$, as in the selected floor package. The denominator
lower bound is also tied to this same bed. At a real momentum $p$, write
\[
 D_i^{\rm bed}(p)=
 \mathsf T\!\left[P_1(-p)^TM_i(p)P_1(p)
       +P_0(p)P_0(-p)^T\right]+s_\mu(p)I
\]
for the denominator matrix in (F13). If
$\mathrm{peb}(p,\mu)$ denotes the positive quadratic weight, the cold
geometry gives
\[
 T^2\leq64\,\mathrm{peb}(p,\mu).
\]
The bed-floor estimate for this matrix gives
\[
 \frac{\min(\epsilon,1)}{2^{10}}\,\mathrm{peb}(p,\mu)\|z\|^2
       \leq \operatorname{Re}\langle z,\mathrm{dais}(p)z\rangle.
\]
Using the denominator matching to identify this matrix with the selected
total denominator, we obtain
\[
 cT^2\|z\|^2\leq
   \operatorname{Re}\langle z,S_i(0)z\rangle,
 \qquad c=\frac{\min(\epsilon,1)}{2^{16}}.
\]
The same cold lower bound and the free reference floor give the ghost
denominator bound with $c=2^{-16}$. For either square matrix $S$, the
coercive estimate implies
$cT^2\|z\|^2\leq\operatorname{Re}\langle z,Sz\rangle
\leq\|z\|\,\|Sz\|$. Hence $\|Sz\|\geq cT^2\|z\|$; in finite
dimension this makes $S$ invertible with
$\|S^{-1}\|_{\mathrm{op}}\leq(cT^2)^{-1}$.

\paragraph{Trajectory prices.}
We derive (F6) from the row contacts and the regulated reference field.
At scale $T_i$, a path estimate with tab budget $L_i$ consists of three
component bounds. The tab derivatives satisfy
$T_i^j\|M_i^{(j)}(0)\|_{\mathrm{op}}\leq C L_i$ for $j=0,1,2$; the
one-leg slot has budget $L_i/T_i$, and the two-leg arch has budget
$L_i/T_i^2$. These are separate components of the path estimate; the
spot and arch bounds do not follow from the tab bound alone. For the interacting row,
the compact derivative caps come from the common strip and selected
fees. Along $x_i+t e_i$, each derivative is a momentum derivative
contracted with $e_i$, whose norm is at most one. If an analytic entry
has radial order $m$, all Taylor tensors of degree below $m$ vanish:
their diagonal polynomial vanishes in every real direction, hence the
polynomial itself is zero. Taylor's integral remainder applied to its
$j$th derivative therefore bounds it by $C\|x_i\|^{m-j}$ for $j<m$;
for $j=m$ use the compact derivative cap. Since $\|x_i\|\leq T_i\leq
r_*$, this gives the path budget $C T_i^m$ for $j\leq m$. Thus the
trajectory transfer with tab budget $L_i$ gives
\[
 \mathsf P_{T_i}(\text{tab};C L_i),\qquad
 \mathsf P_{T_i}(\text{spot};C L_i/T_i),\qquad
 \mathsf P_{T_i}(\text{arch};C L_i/T_i^2).
 \tag{F17}
\]
Thus the interacting transfer uses $L_i=T_i^2$, while the regulated
reference transfer uses $L_i=\mu_i^2$.

For the interacting row, (F14) says the tab is $|p|^2I+O(|p|^3)$;
therefore its path has tab budget $C T_i^2$. For a marked leg, (F15)
says it is the linear vector term plus a quadratic remainder. Along the
path, the linear term has size $O(T_i)$, first derivative $O(1)$, and
zero second derivative. The quadratic remainder has
$T_i^j$-weighted $j$th derivative $O(T_i^2)$ for $j=0,1,2$; since
$T_i\leq r_*$, this is at most $O(T_i)$. Thus the leg has budget
$C T_i$. The
uniform arch cap gives budget $C$ at unit scale. The regulated reference
field has the following three trajectory budgets, which we derive from
its actual construction below:
\[
  C\mu_i^2,\qquad C\frac{\mu_i^2}{T_i},\qquad
  C\frac{\mu_i^2}{T_i^2}.
\]
These weights follow from its construction. Begin with the massive free
film, the free Wilson film plus $\mu_i^2 I$. The free Wilson path has
tab budget $T_i^2$: its tab has a quadratic zero, while its compact
derivatives are uniformly bounded. Its marked leg has a linear first
term and a quadratic remainder, hence budget $T_i$; its two-leg arch
has a uniform unit-scale cap. The added constant mass has tab budget
$T_i^2$ since $\mu_i\leq T_i$, and has zero spot and arch. The massive
film therefore has the three-component path budgets $T_i^2,T_i,1$. Its
scalar tab is the free
dispersion plus $\mu_i^2$; the cold geometric inequality bounds it below
by $T_i^2/64$. Thus its inverse has norm at most $64T_i^{-2}$ and its
inverse path has tab budget $T_i^{-2}$. If $C=A^{-1}$, $B_i$ is the
original one-leg slot, and $D_{ij}$ is the original two-leg slot, the
inverse film's insertion formulas, with $\widehat B_i$ and
$\widehat D_{ij}$ denoting the inverse film's slots, are
\[
 \widehat B_i(p,q)=-C(q+p)B_i(p,q)C(q),
\]
\[
 \widehat D_{ij}(p,q,r)=-\bigl[\widehat B_j(q,r+p)B_i(p,r)
       +\widehat B_i(p,r+q)B_j(q,r)
       +C(r+p+q)D_{ij}(p,q,r)\bigr]C(r).
\]
The one-leg term has size $T_i^{-2}\,T_i\,T_i^{-2}=T_i^{-3}$.
Each of the three two-leg summands has size $T_i^{-4}$: the first two
use $T_i^{-3}\,T_i\,T_i^{-2}$, and the last uses
$T_i^{-2}\,1\,T_i^{-2}$. Thus the inverse film has tab, one-leg, and
two-leg budgets $T_i^{-2}$, $T_i^{-3}$, and $T_i^{-4}$. The reference
construction multiplies this inverse film by $\mu_i^2$, giving budgets
$\mu_i^2/T_i^2$, $\mu_i^2/T_i^3$, and $\mu_i^2/T_i^4$. The same
prices hold for derivatives of orders $j=0,1,2$ along every front or
rear path: apply the product rule to each displayed slot formula and use
the corresponding Week estimate for each differentiated factor. The
shifted momenta have path velocities among $e_i,-e_i$ and their sums,
whose norms are bounded by $2$; this changes only fixed constants. Each
derivative of total order $j$ costs $T_i^{-j}$, which is canceled by
the $T_i^j$ factor in the path price. Since
$\mu_i^2/T_i^2\leq1$, the inverse-reference film has the Week estimate
with tab budget $1$, and consequently spot and arch budgets
$T_i^{-1}$ and $T_i^{-2}$.

The reference insertion is formed from nine copies of this
inverse-reference film and then weighted by $\mu_i^2$. For two films
with tab/spot/arch budgets $(L,L/T,L/T^2)$ and $(M,M/T,M/T^2)$, the
product has tab budget $LM$. Its spot formula has two summands, each
bounded by $LM/T$; its arch formula has four summands, each bounded by
$LM/T^2$. Iterating through the nine factors changes only fixed
numerical constants: the tab, spot, and arch rates remain $1$, $T_i^{-1}$,
and $T_i^{-2}$. The final weight gives $\mu_i^2$,
$\mu_i^2/T_i$, and $\mu_i^2/T_i^2$.
Each directional sandwich in the free reference row has constant
unit-budget factors on both sides, so it preserves these prices.
Summing the four directions changes only the common constant. This
proves the displayed reference budgets from the field construction.

Because $\mu_i\leq T_i$, adding the local and reference tab budgets gives
$C(T_i^2+\mu_i^2)\leq 2C T_i^2$. The total denominator has the same
$T_i^2$ tab budget on both the unshifted and boosted path. The bed floor
proved above gives its inverse norm at $t=0$ at most $C/T_i^2$.
Differentiating the inverse with (F2) gives both inverse prices in
(F6a). The reference arch budget is exactly $C\mu_i^2/T_i^2$, and the
interacting arch has unit-scale budget $C$, proving (F6b). The two total
marked legs have budget at most
$C(T_i+\mu_i^2/T_i)\leq2C T_i$, which proves (F6c). All constants here
are uniform over the selected Hstar family because the fees, compact
caps, regulator bound, floor constant and centered-box radius are uniform.

On the shell, set $T_i=\|x_i\|_2>0$. The unregulated reference
denominator along the path is
\[
 R_i(t)=\|x_i+t e_i\|_2^2 I,
 \qquad R_i(0)=T_i^2I.
\]
Its path budget is $C T_i^2$ and its inverse has price $C/T_i^2$ by
(F2). Let $S_i(t)$ be the total regulated denominator. Its difference
from $R_i$ splits into the local tab minus the free quadratic tab, which
has budget $C T_i^3$ by (F14), and the regulator tab, which has budget
$C\mu_i^2$. Thus
\[
 \mathsf P_{T_i}(S_i-R_i;C(T_i^3+\mu_i^2))
 =\mathsf P_{T_i}(S_i-R_i;C\delta_iT_i^2).
\]
Both inverse paths have price $C/T_i^2$; the identity
\[
 S_i^{-1}-R_i^{-1}=-S_i^{-1}(S_i-R_i)R_i^{-1}
\]
and the product rule therefore give the two inverse differences in (F6d)
with price $C\delta_i/T_i^2$.

For the marked legs, subtract the unregulated free leg. The local marked
row and the free Wilson row have the same linear vector coefficient by
the explicit coefficient calculation above. Their difference has
radial order two. Hence the local marked row minus the halo leg has
price $C T_i^2$ along the shell path. The regulated reference leg has
the one-leg budget $C\mu_i^2/T_i$ from its path spot component. Since
\[
 T_i^2\leq\delta_iT_i,
 \qquad \frac{\mu_i^2}{T_i}\leq\delta_iT_i,
\]
their sum has price $C\delta_iT_i$, proving the leg part of (F6d).
This derivation requires $T_i>0$ and is used only on the shell. At
$x_i=0$ the unregulated reference path is singular, so no halo inverse is
asserted there.

\subsection{Uniform path bounds from the selected records}

For completeness, we describe how the path hypotheses used above are
obtained uniformly for the selected family. Each selected local action
has the same fee bound and positive strip width, and the external vacuum
has its own common fee bound on that strip. Their sum therefore has a
uniform fee bound. Cauchy's integral formula on a smaller common strip
bounds every fixed momentum derivative of each matrix entry uniformly;
for the tab use derivatives through order three, and for the spot and
arch through order two. The momentum maps entering the three components
are fixed linear maps, so the chain rule multiplies these estimates only
by their fixed operator norms. The matrix dimension is fixed, hence the
entrywise bounds imply operator-norm bounds. The additional fixed free
film is a finite polynomial and has uniform derivative bounds on the
centered box. Adding it gives the compact path caps used in (F17).

The selected action's soft hypothesis is componentwise: its unmarked
layer vanishes to radial order three, its one-mark layer to order two,
and its two-mark layer to order one. The local marked insertion uses the
one-mark layer composed with a linear two-momentum map, so it has the
second-order vanishing required for the local contact. The grey vacuum's
marked row has the same linear term as the free Wilson row: the three
plaquette placements give the displayed coefficient
$v_b$, and the fixed free-film term supplies the complementary part.
Their difference has order two. Thus the selected combined row has
exactly the tab and marked-leg contacts used in (F14)--(F15), with
uniform compact caps.

Finally, the reference field is constructed from the same regulator
$\mu$. Its trajectory tab budget is $\mu^2$, and its two other
components carry the factors $T^{-1}$ and $T^{-2}$. The local field has
tab budget $T^2$. Since $\mu\leq T$, their sum has tab budget at most
$2T^2$. The selected House floor and the cold geometric bound give the
total tab a coercive lower bound $cT^2$; the elementary inverse
calculation preceding (F17) then yields the inverse estimates for both
the unshifted and boosted paths. On the shell, the free comparison is
positive because $T=\|x\|_2>0$, and the same uniform strip/cap constants
apply after resetting the scale. This verifies that the family-level
path bounds used in the face proof are common to all selected rows and
all regulators in the stated range.
